% ID: SEL-D0037-OA-E02
% ESTADO: cotejada. Segunda pasada visual y borrador compilado revisados.
% Fuente: Selectividad/matematicas_aplicadas/Ua2MATEMÁTICAS APLICADAS examen.pdf
% SHA-256: a56739d297d4dc220d6898b24cf729dc556401a6bb292b8ca7cfb58c01c31651
% Página: 1; opción A; ejercicio original 2.

La mosca común solamente vive si la temperatura media de su entorno está comprendida entre $4\,{}^\circ\mathrm{C}$ y $36\,{}^\circ\mathrm{C}$. La vida en días, en función de la temperatura media $T$, medida en grados centígrados, viene dada por la función:
\[
V(T)=\frac{-1}{16}\left(T^2-40T+16\right)\qquad T\in[4,36].
\]
\begin{enumerate}[label=\alph*)]
\item \textbf{(1 punto)} Determine la vida máxima que puede alcanzar la mosca común.
\item \textbf{(1 punto)} Calcule la vida mínima e indique la temperatura media a la que se alcanza.
\item \textbf{(0.5 puntos)} Si sabemos que una mosca ha vivido 15 días, \textquestiondown{}a qué temperatura media ha estado el entorno donde ha habitado?
\end{enumerate}
